\section{Introduction}
\label{sec:introduction}

Sparse linear and semidefinite programs are an attractive test case for
quantum optimization.  Their data may be queried or applied without reading a
dense matrix, and an interior-point method reduces optimization to a sequence
of structured linear systems.  A quantum linear-system solver (QLS) can, in
favorable access models, prepare a state proportional to a Newton direction
without first writing every coordinate.  These observations motivate a simple
question:
\begin{quote}
Do quantum interior-point methods beat classical interior-point methods on
sparse problems?
\end{quote}
The word ``sparse'' does not by itself determine the answer.  An input matrix
can be sparse while its normal equations, null-space basis, factorization, or
inverse is dense.  Conversely, a matrix that looks unfavorable through its
condition number can have a clustered spectrum that a classical Krylov method
uses efficiently.  A quantum state proportional to a direction is not yet a
classical iterate, and a preconditioned matrix oracle is not free merely
because the matrix it represents is well conditioned.  The relevant
comparison must therefore name the Newton formulation, the input oracle, the
block-encoding normalization, the right-hand side, the requested output, and
the classical solver that receives the same structure.

This paper develops that comparison at three levels.  At the \emph{geometric}
level, we ask which central-path conditioning is inherent in the optimization
problem and which can be changed by selecting another barrier.  At the
\emph{access-model} level, we ask what a classical or quantum linear solver can
extract from the same spectral structure and what it costs to construct the
operator and right-hand side it receives.  At the \emph{interface} level, we
ask whether the output of one solve can be recovered, certified, and used to
form the next solve.  The resulting picture is neither a blanket quantum
speedup nor a blanket no-go theorem.  It is a set of sharp laws and scoped
frontiers that identify where either conclusion would still have to be
proved.

\subsection{State of play}
\label{sec:introduction-state-of-play}

The first generation of QIPMs inserted a QLS into a classical path-following
framework.  Kerenidis and Prakash treated LPs and SDPs; Casares and
Martin-Delgado developed a predictor--corrector LP method; Augustino et al.
gave conic and semidefinite variants; and the Terlaky-group line developed
inexact feasible, infeasible, orthogonal-subspace, and refinement-based
methods
\cite{kerenidis2020-quantum-interior-point-method-lps,
casares2020-quantum-interior-point-predictor-corrector,
augustino2023-quantum-interior-point-methods-semidefinite,
brandon2023-quantum-interior-point-methods-conic,
augustino2021-an-infeasible-inexact-quantum-interior,
augustino2022-quantum-interior-point-methods-for,
wu2023-inexact-feasible-quantum-interior-point,
mohammadisiahroudi2024-efficient-use-quantum-linear-system,
mohammadisiahroudi2025-inexact-feasible-interior-point-method}.
These methods made the potential role of QLS state preparation precise, but
their ledgers also exposed condition, normalization, tomography, QRAM, and
classical-update costs.  Dalzell et al. emphasize the same end-to-end
accounting issues in their QIPM review and resource analysis
\cite{dalzell2025-quantum-interior-point-methods,
dalzell2023-end-end-resource-analysis-quantum}.

Two newer lines change the architecture rather than merely improving the
linear solver.  The quantum central-path method of Augustino, Leng,
Nannicini, Terlaky, and Wu uses Hamiltonian evolution and potential
evaluation rather than repeated QLS tomography
\cite{augustino2024-quantum-central-path-algorithm-linear}.  The
self-concordant Schr\"odinger-operator program of Gribling et al. likewise
seeks optimization without a Newton condition-number factor
\cite{gribling2025-self-concordant-schrodinger-operators-spectral}.  These
approaches show why a conditioning theorem must be scoped to a named Newton
matrix: a lower bound on a barrier Hessian is not a lower bound on every
quantum optimization architecture.

Apers and Gribling take a different condition-number-free route.  Their tall
LP algorithms use quantum spectral approximation and local-norm estimation,
return an explicit solution, and prove a row-query separation when the number
of constraints is sufficiently larger than the variable dimension
\cite{apers2026-quantum-speedups-linear-programming-interior}.  Within the
QIPM evidence compared here, this is the only proved asymptotic
quantum--classical separation.  It is a row-query separation under a succinct
row oracle, not automatically a wall-clock separation from the best
\(\nnz\)-aware sparse solver.  Its requirements and high powers of the
variable dimension also make it a specific positive regime rather than a
generic consequence of sparsity.

The QLS primitive itself has evolved from HHL through the
Childs--Kothari--Somma precision improvement and QSVT pseudoinversion
\cite{harrow2009-quantum-algorithm-linear-systems-equations,
andrew2017-quantum-algorithm-systems-linear-equations,
andras2019-quantum-singular-value-transformation-beyond}.  Orsucci and Dunjko
proved access-sensitive upper and lower bounds for positive-definite systems.
Mori et al. proved an \(\Omega(\kappa\sqrt{s})\) sparse-access lower bound at
fixed constant error and restated an \(\Omega(\kappa\log(1/\epsilon))\)
constant-sparsity bound from earlier unpublished work of Harrow and Kothari.
Dalzell, Li, and Su introduced truncation and filtering parameters that can be
smaller than a worst-case condition number
\cite{orsucci2021-on-solving-classes-of-positive,
mori2026-sparsity-dependent-complexity-lower-bound,
dalzell2026-faster-quantum-linear-system-solver}.  Somma and de Wolf's related
plain-versus-sum-of-squares separation concerns guided ground-state energy,
not QLS, and we keep that distinction explicit in
\cref{sec:spectra-quantum}
\cite{somma2026-optimal-lower-bound-ground-state-energy}.  These results make
``the QLS cost'' an incomplete phrase: cost depends on whether one has plain
matrix access, factor access, sparse-entry access, or a stronger supplied
oracle, and on how the right-hand side overlaps the hard spectral subspace.

Tomography is the second standard bottleneck.  A path-following method usually
needs a classical direction to update all primal and dual coordinates, while
a QLS naturally returns only a normalized state.  State-preparation-unitary
tomography improves the dependence on readout accuracy, but still scales with
the direction dimension
\cite{mohammadisiahroudi2023-exponentially-more-precise-tomography-for,
joran2023-quantum-tomography-state-preparation-unitaries}.  Avoiding that cost
requires a different output contract and, crucially, a way to generate the
next changing diagonal or Newton operator from the compressed representation.
The interface cannot be removed by renaming a state as an iterate.

Classical sparsity exploitation sets a demanding baseline.  Practical IPMs
choose between normal equations and augmented systems, control fill, reuse
symbolic factorizations and preconditioners, and adapt inner tolerances
\cite{wright1997-primal-dual-interior-point-methods,
jacek2012-interior-point-methods-25-years,
gondzio2023-general-purpose-preconditioning-regularized-interior-point}.
For bounded-treewidth structure, supplied elimination orders and specialized
IPMs can be nearly linear in the problem size
\cite{dong2020-a-nearly-linear-time-algorithm,
gu2022-a-faster-small-treewidth-sdp}.  A sparse-QIPM claim must therefore be
compared with a structure-aware factorization or Krylov method, not with dense
Gaussian elimination by default.

\subsection{What a matched comparison counts}
\label{sec:introduction-ledgers}

The algorithmic literature above cannot be compared by retaining only the
dimension and the sparsity of the original constraint matrix.  Its QIPM lines
use different Newton formulations, state-preparation primitives, refinement
schemes, and outputs.  We use a raw-query/QRAM-write ledger throughout the
paper, counting calls after the oracle is precisely defined and writes needed
to build or refresh QRAM.  We pair it with qualitative end-to-end audits of
construction, reuse, recovery, and output.  Gate and arithmetic work are
populated only for selected results, so these audits are not a uniform
wall-clock ledger.  Neither accounting view substitutes for the other.

\paragraph{From sparse input to the solved operator.}
For an LP, a diagonal scaling of a sparse constraint matrix may produce a
normal equation \(AD^2A^\top\), an augmented KKT matrix, or an
orthogonal-subspace system.  Constant row sparsity of \(A\) need not imply
constant sparsity of any of these objects: one dense column creates a clique
in the normal-equation graph, and a null-space basis can be dense.  The
opposite can also occur, since a carefully chosen basis or augmented
formulation can preserve structure that normal equations lose.  Consequently
our notation distinguishes \(\nnz(A)\), the sparsity \(s_H\) of the actual
system \(H\), and the access description of a factor \(G\) when
\(H=G^\dagger G\); the formal definitions are in
\cref{sec:newton-formulations,sec:prelim-quantum-access}.  A claim that starts
from sparse \(A\) must prove the access properties of the operator it actually
solves.

\paragraph{Condition and normalization.}
A block encoding presents \(H/\alpha_H\), so the inversion scale is governed
by \(\alpha_H\lVert H^{-1}\rVert\), not by the spectral condition number in
isolation.  Rescaling \(H\) changes neither this product nor the underlying
problem.  Multiplying separately encoded factors can multiply their
normalizations, and forming a perfect preconditioned product can require the
same input information that the original solve was meant to discover.  The
factor, congruence, and parity results of \cref{sec:preconditioning} make these
exchanges exact on representative families.  On the classical side, an
effective preconditioner also has a construction and application cost, but it
can often be reused and can exploit entries already stored explicitly.  A
matched comparison prices both constructions rather than granting either side
an ideal inverse oracle.

\paragraph{Right-hand sides and solve accuracy.}
Worst-case QLS complexity is not always the right estimate for a Newton step.
The direction depends on a particular RHS, and truncation can ignore a hard
eigenspace only if the requested output error permits losing its solution
mass.  A residual guarantee, a normalized-state error, and a coordinatewise
direction error are different contracts.  They are connected by
condition-dependent inequalities, so comparing a quantum state-error bound
with a classical residual tolerance without conversion is not meaningful.  The
coupling and spectral-measure analysis in \cref{sec:beyond-kappa} evaluates
these distinctions on exact IPM systems, while the output contracts of
\cref{sec:prelim-output-contracts} state which error is being promised.

\paragraph{State preparation and readout.}
Under normalized operator and RHS access, a QLS call produces a quantum state
proportional to a solution.  A hybrid QIPM that updates a classical iterate
must then estimate a full vector, recover its norm and scale, and ensure that
the resulting step satisfies the outer neighborhood condition.  Repeating the
state preparation for tomography repeats every nonpersistent oracle and RHS
cost.  A method that requests only an observable or compact certificate can
avoid full readout, but then it must show how that object is sufficient for
optimization and for constructing the next iteration.  The distinction among
classical-vector, compact, state, density, coordinate-oracle, and heralded
outputs in \cref{sec:prelim-output-contracts} is designed to prevent a proof
for one interface from being silently reused for another.

\paragraph{The classical comparator.}
For a sparse direct method, the per-step cost is controlled by realized fill,
numerical factorization, and triangular solves, with a symbolic ordering
normally reused.  Multiple predictor or corrector RHSs can reuse the same
factorization.  For a Krylov method, the relevant costs are matrix--vector and
preconditioner applications times a spectrum- and tolerance-dependent
iteration count.  Clustered eigenvalues can make that count far smaller than
a worst-case \(\sqrt\kappa\) estimate, as
\cref{prop:two-cluster-cg} demonstrates.  Specialized graph, treewidth,
block-angular, and low-rank structure can improve the baseline further.  We
therefore do not compare a quantum query count directly with dense cubic
arithmetic or treat the act of writing a dense classical direction as the
only classical cost.

\paragraph{Setup, refresh, and trajectory reuse.}
One-step costs do not automatically multiply by the IPM iteration count.  A
fixed sparse input can be read once, a symbolic factorization can persist,
and a lazy oracle can combine stored data with a changing scalar path
parameter.  Conversely, a QRAM or block encoding of changing primal and slack
diagonals may require per-step writes, and tomography may be repeated at every
step.  The distinction is formalized by the amortized accounting contract in
\cref{def:amortized-accounting}.  Our online direct-sum lower bound applies
only when new oracle information is genuinely fresh
\cref{thm:prec-online-direct-sum}; the caching-ceiling argument in
\cref{sec:preconditioning-lazy} shows why no unqualified
\(T\)-times-one-step claim follows for a static LP.

\paragraph{Query separation versus time separation.}
Raw coefficient, sparse-entry, row, block, QRAM, and preconditioned-product
queries reveal different amounts of information.  A sublinear row-query
algorithm can be a genuine oracle separation even when generating a queried
row or performing local arithmetic prevents a time advantage.  Conversely,
an explicitly stored sparse matrix may have no interesting input-query cost,
while its repeated numerical solves remain expensive.  The Apers--Gribling
result is positive in its declared row-query model; our parity and holonomy
lower bounds are negative in their declared raw local-query models.  Neither
statement changes models for the sake of the comparison.  This discipline is
why the later theorems often conclude with a setup--solve--recovery inequality
rather than with a single condition-number exponent.

\subsection{Answer in outline}
\label{sec:introduction-answer}

Our first answer is that central-path conditioning is geometry before it is
barrier design.  For a compact feasible slice, the diameter of an objective
sublevel set and the length of its longest chord control the extreme
eigenvalues of every self-concordant barrier Hessian.  The sublevel chord law
in \cref{thm:sublevel-chord-law} gives a condition-number floor at a fixed
objective gap, and \cref{thm:canonical-achievability} shows that the canonical
log and log-det barriers attain that floor up to explicit instance constants.
The conclusion is minimax and gap matched: changing the barrier cannot in
general remove the divergence.  It does not say that every optimum is ill
conditioned, nor does it cover a method that avoids barrier-Newton solves.

Our second answer is that a large condition number does not determine the
relative classical and quantum costs.  Degenerate LP log-barrier Hessians
split into two intervals with bounded internal ratios
\cref{prop:two-cluster-hessian}.  A classical residual polynomial can exploit
the empty interval, and \cref{prop:two-cluster-cg} gives polylogarithmic
iteration dependence on the inter-cluster separation.  Quantum polynomial
transformations must also remain bounded on the domain imposed by their
access model.  Under plain block-encoding access the worst-case QLS cost is
\(\Thetatilde(\kappa)\), even for two-point clusters when
\(N_{\rm QLS}\gtrsim\kappa\).  A norm-tight factor presentation
unconditionally has a \(\Thetatilde(\sqrt\kappa)\) polynomial route in its
model; only the corresponding end-to-end solve needs the overlap hypotheses
in \cref{thm:quantum-cluster-obstruction}.  On-path right-hand sides can be
much easier still: \cref{prop:newton-rhs-alignment,cor:top-window-solve} show
that the exact central-path right-hand side has \(O(\mu^2)\) weak-cluster
fraction, so a top-window solve meets the fixed relative-residual contract
once the spectral gap is resolved.

This classical--quantum contrast is not paradoxical.  CG is charged for
matrix--vector products and can place roots in spectral gaps.  QSVT is charged
for a bounded transformation of a supplied normalized operator, and the
normalization and symmetric-domain constraints remain part of the problem.
The comparison reverses the naive view that a large \(\kappa\) is always bad
for the classical method and automatically favorable to a quantum method.
It also motivates the instance-selection laws of
\cref{sec:beyond-kappa}: one should measure right-hand-side coupling and
spectral mass, not select a solver from \(\kappa\) alone.

Our third answer is that the decisive lower bounds often live outside the
linear solve.  The LP constructions of \cref{sec:lp-hardness} can make the
reduced Newton geometry condition one while hiding a parity-dependent affine
translation or original-coordinate direction.  The SDP constructions of
\cref{sec:sdp-hardness} separate intrinsic value information, Schur solves,
KKT access, and recovery.  The preconditioning results of
\cref{sec:preconditioning} show product laws: improving spectral condition
can worsen factor normalization or physical-state sensitivity, and a state
from one solve need not provide a coordinate oracle for the next.  The right
unit of analysis is therefore a complete setup--load--solve--recover--output
ledger.

Condensation supplies a complementary positive and negative example.  A
shared trace budget drives winner-take-all mass onto the lowest-cost PSD
block.  That mass simultaneously controls reduced-Hessian conditioning,
trace-distance visibility, and the query cost of state preparation.  The
transition has matching lower and upper query bounds on a bounded-incidence
holonomy SDP, while small-success search gives a tight accuracy curve on its
proved domain; the \(\delta=o(k/G)\) sliver and the internal-state
\(\eta_{\rm int}\)-channel remain open
\cref{sec:state-conversion-open}.  Thus output hardness is not inferred from
condition alone: it is proved through a visible, decodable output contract.

Finally, our no-go results are frontier theorems, not universal
impossibilities.  Convex mixtures, cut resistance, bounded-arity Walsh rows,
symmetry work identities, scalar congruences, and state-interface products
rule out the natural gadget classes to which their hypotheses apply.  The
section records the five escape routes in \cref{sec:frontier-synthesis}.  The
positive modules of
\cref{sec:positive-modules} inhabit some of those routes, but remain
conditional and per-step: they require stable geometry, explicit access,
compressed output, or a low-dimensional border.  No section silently upgrades
one such module to an end-to-end sparse-QIPM advantage.

\subsection{Contributions}
\label{sec:introduction-contributions}

The technical contributions follow the paper's order.  Each item names the
result that supports its headline and the boundary under which it should be
read.

\begin{enumerate}[leftmargin=*]
 \item \textbf{Universal conditioning geometry (\cref{sec:conditioning}).}
 We prove minimum-eigenvalue rigidity and the sublevel chord law for every
 self-concordant barrier
 \cref{thm:lambda-min-rigidity,thm:sublevel-chord-law}, sharpen it to a
 width-by-width spectral sandwich \cref{thm:width-spectrum-rigidity}, and
 show canonical LP and product-SDP barriers attain the resulting floor up to
 an explicit gap-matched factor \cref{thm:canonical-achievability}.  The
 conditioning dictionary includes positive-dimensional optimal faces,
 degenerate vertices, curved optima, singularity-degree exponents, and
 approximate centrality
 \cref{cor:conditioning-regimes,cor:fractional-conditioning-laws,cor:near-degeneracy-plateau,cor:approximate-centrality-transfer}.

 \item \textbf{Clustered spectra and solver separation
 (\cref{sec:spectra}).}  We prove the two-cluster LP tail
 \cref{prop:two-cluster-hessian}, the cluster-sensitive CG bound
 \cref{prop:two-cluster-cg}, and the plain-versus-factor quantum polynomial
 frontier \cref{thm:quantum-cluster-obstruction}.  The exact central-path
 right-hand side obeys a weak-mode cancellation law, yielding a
 condition-number-free top-window solve under the fixed relative-residual
 contract once the spectral gap is resolved
 \cref{prop:newton-rhs-alignment,cor:top-window-solve}.  The generic
 \(\Thetatilde(\sqrt\kappa)\) factor-access QLS interpretation is explicitly
 rejected: the theorem gives a polynomial metric and an overlap-qualified
 solve, not such an unrestricted algorithm.

 \item \textbf{Instance laws beyond \(\kappa\)
 (\cref{sec:beyond-kappa}).}  Exact block inversion identifies when DLS
 filtering improves on worst-case conditioning
 \cref{thm:dls-coupling}; a truncation transition separates solution-state
 and residual contracts \cref{prop:dls-truncation-transition}; and the
 effective-spectral-dimension theorem places truncation and filtering
 thresholds at dimensions four and eight
 \cref{thm:spectral-dimension-laws}.  Stable active-range equilibration
 preserves the Newton state for \(f\in\operatorname{range}Q\) while removing
 the divergent tail when its projector is supplied
 \cref{thm:stable-split-equilibration}.

 \item \textbf{Block log-det condensation
 (\cref{sec:condensation}).}  We derive the exact capacity transition and
 ground-multiplicity-dependent conditioning phases
 \cref{thm:condensation-threshold,thm:cond-conditioning-phases}, prove
 finite mass--conditioning laws \cref{thm:cond-mass-conditioning}, and link
 them to trace-distance visibility and raw-query complexity
 \cref{thm:cond-visibility,thm:cond-unique-or}.  The transition extends to
 sublinear multiwinner sets and to a stated approximate-centrality model
 \cref{thm:cond-multiwinner,thm:cond-robust}.  These are statements about the
 equality-reduced Frobenius Hessian and trace-normalized output, not every
 preconditioned KKT representation.

 \item \textbf{A sparse winner-take-all SDP
 (\cref{sec:winner-take-all}).}  A bounded-incidence qutrit-holonomy family
 realizes the condensation law with tight \(\Theta(N\sqrt G)\) quantum and
 \(\Theta(NG)\) randomized value-query costs, a condition-one public root
 start, robust winner recovery \cref{thm:wta-robust-optimizer}, and a matching
 search-then-solve compact-output algorithm
 \cref{thm:wta-search-upper}.  The formulation tradeoff in
 \cref{sec:wta-scope} keeps the bounded-incidence claim separate from an
 isometric ambient conditioning claim.

 \item \textbf{Accuracy curves for hidden-block state conversion
 (\cref{sec:state-conversion}).}  A verified-sample decoder gives a
 mixed-output lower bound \cref{thm:sc-m1}; detuned rejection preparation
 supplies the upper bound \cref{thm:sc-m2}; and parity collapse yields the
 tight nontrivial conversion curve
 \cref{thm:sc-small-success,cor:sc-s2}.  The batch results separate
 nonamortizable vector materialization from scalar approximate counting
 \cref{thm:sc-batch-vector,thm:sc-batch-scalar}.  The theorem domains retain
 the loser-state discrepancy, the small-success sliver, and logarithmic
 fixed-point costs rather than extrapolating through them.

 \item \textbf{LP loading and recovery hardness
 (\cref{sec:lp-hardness}).}  The gain--plateau family gives a linear robust
 primal-state lower bound \cref{thm:lp-gain-plateau}.  Prefix rigidification
 strengthens this to condition-one reduced Newton geometry with linear
 original-coordinate recovery cost \cref{thm:lp-prefix-capstone}.  Complete
 KKT and Schur-preconditioned variants isolate public-right-hand-side,
 residual, bounded-data, and scale tradeoffs
 \cref{thm:lp-full-kkt,thm:lp-preconditioned-kkt,thm:lp-bounded-exact,thm:lp-unit-threshold,thm:lp-dual-homogeneous-full}.
 None is presented as a condition-one QLS
 lower bound with free transformed oracles.

 \item \textbf{Gadget-class frontiers
 (\cref{sec:gadget-frontiers}).}  A single-flip convex-mixture theorem gives
 the graph-free default obstruction \cref{thm:frontier-single-flip}.
 Cut--resistance and Walsh-row theorems quantify the volume or scale needed
 by signed-copy and bounded-arity constructions
 \cref{thm:frontier-cut-resistance,thm:frontier-walsh}; gauge-flow, work, and
 Fourier identities locate the multiplier or input-incidence cost
 \cref{thm:frontier-gauge-flow,thm:frontier-work,thm:frontier-fourier-source}.
 The escape routes in
 \cref{sec:frontier-synthesis} are part of the result.

 \item \textbf{SDP-native loading and recovery hardness
 (\cref{sec:sdp-hardness}).}  An intrinsic trace-normalized qutrit family has
 parity-hard value and a public, condition-one reduced Newton step whose
 primal and complete-KKT solution states are parity hard
 \cref{thm:sdp-intrinsic,thm:sdp-newton-state}.  Non-SOCP Schur families give
 an exact loading/accuracy frontier
 \cref{thm:sdp-schur,thm:sdp-accuracy-frontier}, and tree and group holonomy
 lift the mechanism to matrix-valued words
 \cref{thm:sdp-tree-loading,thm:sdp-holonomy-spectrum,thm:sdp-s3-hardness}.
 The proved query exponent for the \(S_3\) family is
 still obtained on an abelian restriction; intrinsically nonabelian word
 hardness remains open.

 \item \textbf{Preconditioning and state-interface products
 (\cref{sec:preconditioning}).}  Whitening and reusable inverse access obey
 normalization tradeoffs
 \cref{thm:prec-whitening-frontier,thm:prec-inverse-access}; scalar
 congruence trades transformed condition against physical recovery
 \cref{thm:prec-congruence-frontier,thm:prec-general-scalar-frontier}; and
 the parity ledger charges setup, RHS preparation, solve, and recovery
 \cref{thm:prec-factor-rhs-hardness,thm:prec-parity-ledger}.  State-to-next-solve
 and lazy-oracle theorems distinguish a one-use state from a reusable
 coordinate service
 \cref{thm:prec-state-diagonal,thm:prec-state-next-solve,thm:prec-lazy-lookup,thm:prec-reusable-endpoints}.

 \item \textbf{Conditional positive modules
 (\cref{sec:positive-modules}).}  We prove a residual-certified pathwise
 refresh bound \cref{thm:positive-lazy-refresh}, a constructive stable
 active-range equilibration and exact face repair
 \cref{thm:positive-partition-free-equilibration,thm:positive-face-repair}, an
 implicit slack oracle and conditional
 compressed dual QIPM
 \cref{thm:positive-implicit-slack,thm:positive-compressed-qipm}, and a
 block-angular border solver whose \(\sqrt N\) block-count exponent is
 optimal in the block-oracle model, as proved in
 \cref{sec:positive-block-angular}.  Every theorem keeps its tail, gap,
 forcing, normalization, and output assumptions visible.

 \item \textbf{Reusable structural tools
 (\cref{sec:structural-tools}).}  Involution and compact-group work
 identities \cref{thm:structural-involution-work,thm:structural-group-work},
 a noncommutative central-mixture sandwich
 \cref{thm:structural-sdp-kantorovich}, and rank and support barriers
 \cref{thm:structural-rank-progress,thm:structural-support-barrier} explain
 several no-go mechanisms abstractly.  Sparse OSS and treewidth-two KKT
 embeddings preserve known QLS restrictions
 \cref{thm:structural-oss-embedding,thm:structural-treewidth-two}, while
 exact-treewidth value hardness and a central-oracle light cone distinguish
 query, output, and physical-locality models
 \cref{thm:structural-exact-treewidth,thm:structural-light-cone}.

 \item \textbf{Corrections and dead ends (\cref{sec:corrections}).}  We show
 that the global potential supremum used in the QCPM simulator ledger cannot
 scale as asserted \cref{prop:corr-qcpm-supremum}, correct the independent
 clock/time-dilation error \cref{eq:corr-qcpm-clock}, derive the corrected norm
 scale and its speed tradeoff
 \cref{thm:corr-qcpm-norm,cor:corr-qcpm-tradeoff}, and give explicit
 densification and hidden-start counterexamples
 \cref{thm:corr-path-simplex,thm:corr-hidden-sign}.  The recorded dead ends
 prevent first-inverse norms, sparse metrics, exact Schur products, or
 one-step bounds from being reused as claims they do not prove.
\end{enumerate}

\subsection{Related work and positioning}
\label{sec:introduction-related-work}

\paragraph{Quantum interior-point methods.}
The QIPM literature spans QLS-based hybrid path following, iterative
refinement, orthogonal-subspace formulations, condition-number-free sampling,
and Hamiltonian central paths.  Representative lines include
Kerenidis--Prakash, Casares--Martin-Delgado, Augustino et al.,
Mohammadisiahroudi et al.,
Wu et al., Apers--Gribling, and the reviews and resource ledgers of Dalzell et
al.
\cite{kerenidis2020-quantum-interior-point-method-lps,
casares2020-quantum-interior-point-predictor-corrector,
augustino2023-quantum-interior-point-methods-semidefinite,
mohammadisiahroudi2025-inexact-feasible-interior-point-method,
wu2023-inexact-feasible-quantum-interior-point,
apers2026-quantum-speedups-linear-programming-interior,
dalzell2025-quantum-interior-point-methods}.  The QCPM and
self-concordant-Schr\"odinger lines avoid the usual repeated-QLS architecture
\cite{augustino2024-quantum-central-path-algorithm-linear,
gribling2025-self-concordant-schrodinger-operators-spectral}.  We do not
collapse these algorithms into a single complexity table: each uses a
different access, condition, precision, and output contract.  The per-section
novelty paragraphs compare the exact theorem at hand with its nearest line of
work.

\paragraph{Quantum linear systems.}
HHL introduced sparse QLS state preparation, CKS improved the precision
dependence, and QSVT placed inversion in a general block-encoding calculus
\cite{harrow2009-quantum-algorithm-linear-systems-equations,
andrew2017-quantum-algorithm-systems-linear-equations,
andras2019-quantum-singular-value-transformation-beyond}.  Orsucci--Dunjko's
positive-definite lower bounds ground our plain-access clustered comparison.
Mori et al. prove \(\Omega(\kappa\sqrt{s})\) at fixed constant error and
restate an \(\Omega(\kappa\log(1/\epsilon))\) constant-sparsity bound from
earlier unpublished work of Harrow and Kothari.  DLS provide the
beyond-condition-number parameters evaluated in
\cref{sec:beyond-kappa}
\cite{orsucci2021-on-solving-classes-of-positive,
mori2026-sparsity-dependent-complexity-lower-bound,
dalzell2026-faster-quantum-linear-system-solver}.  Somma--de Wolf provide an
analogous plain-versus-sum-of-squares resolution separation in a guided
Hamiltonian problem, not a QLS theorem
\cite{somma2026-optimal-lower-bound-ground-state-energy}.  This access-model
discipline is essential to \cref{sec:spectra,sec:preconditioning}.

\paragraph{Quantum optimization lower bounds.}
Van Apeldoorn et al. establish general quantum SDP upper and lower bounds and
Boolean-in-optimum constructions; Brand\~ao--Svore initiated the quantum SDP
solver line; and Apers--Gribling's Theorem~8.4 gives a sparse-LP value lower
bound already having a linear exponent in balanced dimensions
\cite{apeldoorn2020-quantum-sdp-solvers,
fernando2017-quantum-speed-ups-solving-semidefinite,
apers2026-quantum-speedups-linear-programming-interior}.  Our parity and
holonomy families do not claim parity, adversary composition, or the linear
exponent as new.  Their contribution is the conjunction with exact central
paths, condition-one or bounded-condition formulations, robust output
contracts, and complete loading/recovery ledgers described in
\cref{sec:winner-take-all,sec:lp-hardness,sec:sdp-hardness}.

\paragraph{Classical interior-point structure.}
Wright's work on logarithmic-barrier Hessians and numerical error, the
standard primal--dual theory, and the Forsgren--Gill--Wright survey provide
the conditioning and implementation background
\cite{wright1994-properties-hessian-logarithmic-barrier,
wright1998-ill-conditioning-computational-error-interior,
wright1997-primal-dual-interior-point-methods,
forsgren2002-interior-methods-nonlinear-optimization}.  Gondzio and
collaborators document preconditioning, structure exploitation, and practical
IPM evolution
\cite{jacek2012-interior-point-methods-25-years,
gondzio2023-general-purpose-preconditioning-regularized-interior-point}.
Small-treewidth LP and SDP algorithms show why bounded graph width is a
classical algorithmic asset rather than merely a sparse-oracle promise
\cite{dong2020-a-nearly-linear-time-algorithm,
gu2022-a-faster-small-treewidth-sdp}.  The sublevel chord law differs from the
classical log-barrier conditioning analyses by quantifying over all
self-concordant barriers, while \cref{sec:spectra} uses rather than ignores
the classical cluster-sensitive Krylov advantage.

\paragraph{Positioning discipline.}
Many ingredients used below are established: self-concordant containment,
weak-sharp and singularity-degree error bounds, matrix-Burg resolvents,
Rayleigh--Jeans condensation, parity and search, group synchronization,
Kantorovich inequalities, effective resistance, and sparse elimination.  We
credit them where each enters.  Because the nearest prior work changes from
one theorem package to the next, this introduction does not restate every
collision or priority qualification.  The detailed novelty and scope
statements in
\cref{sec:conditioning-overview,sec:condensation-discussion,sec:wta-scope,sec:sdp-synthesis,sec:preconditioning-synthesis,sec:positive-synthesis,sec:structural-synthesis}
are part of the claims.

\subsection{Conventions, verification, and corrections}
\label{sec:introduction-verification}

The distinctions formalized in \cref{sec:preliminaries} are standing
conventions, not optional caveats.  In particular, input sparsity differs from
Newton-system sparsity; \(\kappa_H\) differs from block-encoding normalization;
reduced and ambient KKT conditioning are different quantities; and value,
classical-vector, quantum-state, coordinate-oracle, and heralded outputs are
different contracts.  Query bounds count the setup and transformations named
by their access model.  An uncharged stronger oracle is treated as a stronger
problem statement, not as an implementation of raw sparse access.

The verification process changed several statements inherited from
preliminary research notes and identified two independent issues in a public
preprint.  The corrections relevant
to reading the paper are:
\begin{enumerate}[leftmargin=*]
 \item \cref{sec:corrections-qcpm} corrects the QCPM v2 simulator-norm
 estimate and its independent clock/time-dilation error: a global supremum on
 a fixed truncation box is not controlled by concentration near the
 instantaneous center, and the time-varying kinetic coefficient requires the
 corrected clock in \cref{eq:corr-qcpm-clock}.
 \Cref{thm:corr-qcpm-norm,cor:corr-qcpm-tradeoff} give the resulting scale.
 This is a correction to that derivation, not a general impossibility result
 for Hamiltonian optimization.
 \item \cref{sec:spectra-quantum} corrects an overbroad factor-access
 inference in our earlier analysis.  A square-root polynomial metric, or a
 two-stage solve with explicit overlap, does not imply a generic
 \(\Thetatilde(\sqrt\kappa)\) factor-access QLS; the Orsucci--Dunjko instance
 rules out that unrestricted interpretation.
 \item \cref{sec:spectra-two-cluster} replaces a coarse three-scale picture
 of the degree-two SDP witness by the verified four-window spectrum.  The
 correction changes the intermediate spectral description, not the
 condition-number exponent.
 \item \cref{sec:effective-spectral-dimension} anchors effective spectral
 dimension at the smallest visible eigenvalue.  This repairs the earlier
 shell formulation and is needed for the dimension-four and dimension-eight
 phase laws in \cref{thm:spectral-dimension-laws}.
 \item \cref{sec:state-conversion-model,sec:state-conversion-small-success,sec:state-conversion-open}
 replace the collision shortcut by a verified-sample decoder, restrict the
 classical linear-in-success law to its actual domain, keep the
 very-small-success sliver open, and attribute the \(k\)-marked small-success
 search bound to Dohotaru--H\o yer rather than to Zalka's one-marked theorem
 alone.
 \item \cref{sec:positive-face-repair} corrects the earlier claim that two
 convergent leakage ledgers alone bound projective variation.  Neither the
 ledgers nor an additional $O(\mu_t)$ tail bound suffices: an
 oscillating-face-displacement counterexample has convergent ledgers, the tail
 bound, and linearly growing realized variation, so finite realized variation
 stays a separate hypothesis.
\end{enumerate}
The list is deliberately narrow.  It records changes that affect the theorem
statements or their interpretation, rather than presenting ordinary proof
repairs as additional results.  Further failed shortcuts and their valid
replacements are collected in \cref{sec:corrections-dead-ends}.

\subsection{Organization and reading paths}
\label{sec:introduction-organization}

The conventions are fixed in \cref{sec:preliminaries}; thereafter we use the
section pointers rather than redefining those contracts.  Conditioning
geometry, clustered spectra, and beyond-condition-number instance laws occupy
\cref{sec:conditioning,sec:spectra,sec:beyond-kappa}.  Condensation, its sparse
winner-take-all realization, and the associated state-conversion accuracy
curve occupy \cref{sec:condensation,sec:winner-take-all,sec:state-conversion}.
The LP and SDP loading/recovery lower bounds, gadget frontiers, and
preconditioning/interface obstructions are in
\cref{sec:lp-hardness,sec:gadget-frontiers,sec:sdp-hardness,sec:preconditioning}.
The conditional positive modules are in
\cref{sec:positive-modules}, structural tools in \cref{sec:structural-tools},
corrections and dead ends in \cref{sec:corrections}, and the combined verdict
and open frontier in \cref{sec:discussion}.

Three shorter paths may be useful.  A reader interested only in conditioning
can read
\cref{sec:preliminaries,sec:conditioning,sec:spectra,sec:beyond-kappa}.  A
reader interested only in lower bounds can read
\cref{sec:preliminaries,sec:condensation,sec:winner-take-all,sec:state-conversion,sec:lp-hardness,sec:gadget-frontiers,sec:sdp-hardness,sec:preconditioning}.
A reader interested only in algorithmic
modules can read
\cref{sec:preliminaries,sec:beyond-kappa,sec:positive-modules}.  In every
path, the output and access definitions of \cref{sec:preliminaries} remain
part of the theorem statements.
